\section{Hauptkapitel}
\label{chap:kern}

Das System implementiert den Kernablauf eines Predictionsmarkts. Eine Adresse legt
einen Markt mit Frage, Auflösungskontext, Laufzeit und Ausgängen an. Andere
Adressen beziehen Token über das Faucet, genehmigen den Marktvertrag und staken
auf einen Ausgang. Solange die Frist läuft, sammeln sich Einsätze in den Pools.
Nach Fristende setzt die Admin-Adresse den gewinnenden Ausgang. Gewinner fordern
ihren Anteil in einer eigenen Transaktion an.

Die Benutzeroberfläche reduziert neue Märkte auf \emph{YES} und \emph{NO} für
die Demonstration. Der Vertrag unterstützt formal mehrere Ausgänge. Die
Implementierung bleibt dabei transparent und nachvollziehbar.

Die Implementierung beruht auf drei Entscheidungen. Das Faucet-Modell senkt die
Einstiegshürde, da keine reale Wertposition nötig ist. Der ERC-20-Ablauf mit
\texttt{approve} und \texttt{transferFrom} führt zu einem zweistufigen Zahlungs-
fluss. Die Admin-basierte Auflösung schafft einen zentralen Vertrauenspunkt.

Gas- und Benutzungsaspekte sind reduziert. Die Anwendung zeigt nur die für den
Kernpfad notwendigen Transaktionen und vermeidet zusätzliche on-chain Schleifen
oder Sammelauszahlungen. Das Ergebnis ist ein schmaler, leicht verständlicher
Funktionsumfang.
